% !TeX program = pdflatex
\documentclass[a4paper,14pt]{extarticle}

\usepackage[T2A]{fontenc}
\usepackage[utf8]{inputenc}
\usepackage[russian]{babel}
\usepackage{mathptmx}
\usepackage[a4paper,left=25mm,right=15mm,top=20mm,bottom=20mm]{geometry}
\usepackage{amsmath,amssymb}
\usepackage{array,booktabs,tabularx,longtable}
\usepackage{float}
\usepackage{caption}
\usepackage{enumitem}
\usepackage{xcolor}
\usepackage{tikz}
\usetikzlibrary{shapes.geometric,arrows.meta,positioning,calc}
\usepackage[hidelinks]{hyperref}

\captionsetup[table]{skip=4pt,font=small}
\captionsetup[figure]{skip=4pt,font=small}
\setlist[itemize]{leftmargin=1.6em,itemsep=2pt,topsep=2pt}
\setlist[enumerate]{leftmargin=1.8em,itemsep=2pt,topsep=2pt}

% ---------- данные для титульного листа ----------
\newcommand{\university}{МИНОБРНАУКИ РОССИИ\\[2pt]
Федеральное государственное бюджетное образовательное учреждение\\
высшего образования\\[2pt]
<<\underline{\hspace{9cm}}>>}
\newcommand{\department}{Кафедра \underline{\hspace{7cm}}}
\newcommand{\discipline}{Теория автоматов}
\newcommand{\studentgroup}{\underline{\hspace{3.2cm}}}
\newcommand{\studentname}{\underline{\hspace{6cm}}}
\newcommand{\teacher}{\underline{\hspace{6cm}}}
\newcommand{\labtitle}{Турнир ботов <<камень --- ножницы --- бумага>>:\\[2pt]
конечный автомат с поиском паттерна в истории}
\newcommand{\labnumber}{2}
\newcommand{\cityyear}{Москва, 2026}

\begin{document}

% =====================================================================  ТИТУЛ
\begin{titlepage}
\centering
{\small\university\par}
\vspace{1.5cm}
{\large\department\par}
\vspace{2.5cm}
{\Large\bfseries ОТЧЁТ\\[4pt] по лабораторной работе № \labnumber\par}
\vspace{0.6cm}
{\large по дисциплине <<\discipline>>\par}
\vspace{1.5cm}
{\Large\bfseries\labtitle\par}
\vfill
\begin{flushright}
\begin{tabular}{@{}l@{}}
Выполнил:\\
студент группы \studentgroup\\
\studentname\\[10pt]
Проверил:\\
\teacher\\
\end{tabular}
\end{flushright}
\vfill
{\cityyear\par}
\end{titlepage}

\tableofcontents
\newpage

% =====================================================================  1
\section{Цель работы}

Цель работы --- спроектировать конечный автомат, управляющий ботом в турнире
<<камень --- ножницы --- бумага>>, оформить его в виде схемы алгоритма по
ГОСТ~19.701--90, построить таблицы, задающие работу автомата, и проверить его
на примере турнирного матча.

Автомат должен быть \emph{детерминированным}: генератор случайных чисел не
используется, каждый ход однозначно выводится из истории игры.

% =====================================================================  2
\section{Обозначения и начальные установки}

В работе используются только перечисленные ниже переменные. Никакие другие
величины в схеме и таблицах не вводятся: число раундов в матче ($5$) и глубина
паттерна ($2$) --- это константы, а не переменные.

\begin{table}[H]
\centering
\caption{Переменные алгоритма}
\label{tab:vars}
\begin{tabular}{|c|p{11.4cm}|}
\hline
\textbf{Обозн.} & \textbf{Смысл}\\
\hline
МХ  & мой ход (выход автомата)\\
\hline
ИСТ & история игры --- таблица~4 (список записей о сыгранных раундах)\\
\hline
НР  & номер раунда\\
\hline
ХС  & ход соперника (приходит от сервера)\\
\hline
СХС & следующий ход соперника (прогноз, получаемый из таблицы~5)\\
\hline
РЗ  & результат раунда: В --- победа, Н --- ничья, П --- поражение\\
\hline
КХ  & контрход --- ход, который бот готовит на следующий раунд\\
\hline
СП  & счётчик поражений\\
\hline
СН  & счётчик ничьих\\
\hline
С   & стратегия: $0$ --- стратегия~1, $1$ --- стратегия~2\\
\hline
\end{tabular}
\end{table}

Начальные установки (блок <<Н.У.>> в схеме) задаются строго так:
\begin{equation}
\boxed{\;\text{НР}:=0;\quad \text{СП}:=0;\quad \text{СН}:=0;\quad
\text{С}:=0;\quad \text{ИСТ}:=\varnothing\;}
\label{eq:init}
\end{equation}

Ходы кодируются буквами К (камень), Н (ножницы), Б (бумага). Отношение
<<бьёт>> задаётся функцией $b$:
\begin{equation}
b(\text{К})=\text{Б},\qquad b(\text{Н})=\text{К},\qquad b(\text{Б})=\text{Н},
\label{eq:b}
\end{equation}
то есть $b(x)$ --- ход, которым можно побить ход $x$.

% =====================================================================  3
\section{Постановка задачи}

Матч состоит из пяти раундов. В каждом раунде бот и соперник одновременно
называют ход; сервер сравнивает их и сообщает результат. Требуется, чтобы бот:

\begin{enumerate}
\item не использовал случайность и не зависел от номера раунда как от
      <<расписания>> ходов;
\item учитывал поведение конкретного соперника, то есть \emph{учился} на
      истории;
\item оставался конечным автоматом --- с конечной памятью и конечным числом
      состояний.
\end{enumerate}

% =====================================================================  4
\section{Метод: две стратегии}

\subsection{Стратегия 1 --- реактивная}

Простейший разумный ответ: побить последний ход соперника. Если соперник
только что сыграл $x$, бот играет $b(x)$. Это задаёт \textbf{таблицу~1}
(контрход стратегии~1).

\subsection{Стратегия 2 --- проактивная (по паттерну)}

Стратегия~1 всегда отстаёт на один ход: она реагирует на уже случившееся.
Стратегия~2 пытается \emph{предсказать} следующий ход соперника.

Гипотеза: соперник --- не генератор случайных чисел, а программа или человек,
и в его последовательности ходов есть повторяющиеся фрагменты (паттерны).
Если последние $k$ ходов соперника уже встречались раньше, то разумно ожидать,
что и продолжение будет таким же, как тогда.

Алгоритм прогноза:
\begin{enumerate}
\item взять последние $k$ ходов соперника из ИСТ --- это паттерн;
\item найти в ИСТ более раннее (не совпадающее с самим паттерном) вхождение
      такого же паттерна;
\item посмотреть, какой ход соперник сыграл сразу после него, --- это и есть
      прогноз СХС.
\end{enumerate}

Поиск выполняется с перебором глубины: сначала $k=2$ (пары), и только если пар
в истории ещё не повторялось --- $k=1$ (отдельные ходы). Такой поиск задаёт
\textbf{таблицу~5}.

Получив прогноз СХС, бот играет контрход $b(\text{СХС})$ --- это
\textbf{таблица~2}.

\subsection{Почему таблицы 1 и 2 совпадают}

В таблицах~1 и~2 стоит одна и та же функция $b(\cdot)$: и там, и там бот
выбирает ход, побивающий некоторый ход соперника. Различие --- в
\emph{аргументе}:
\[
\text{табл.~1: } \text{КХ}=b(\text{ХС})\ \text{(фактический ход)},\qquad
\text{табл.~2: } \text{КХ}=b(\text{СХС})\ \text{(прогноз)}.
\]
Таблицы оставлены раздельными, потому что это две разные стратегии с разными
источниками данных; совпадение функций --- свойство игры, а не ошибка.

\subsection{Правило переключения}

Стратегия~1 работает, пока бот не накопил двух неудачных раундов. Как только
\begin{equation}
\text{СП}+\text{СН}\ge 2,
\label{eq:switch}
\end{equation}
включается стратегия~2, и переключение \emph{необратимо}: $\text{С}:=1$.
Смысл порога: двух ничьих или поражений достаточно, чтобы признать, что
реактивная тактика не работает, и перейти к попытке прогноза.

Если стратегия~2 включена, но паттерн не найден (история ещё слишком короткая
или не содержит повторов), бот не остаётся без хода --- он играет по
таблице~1. Это резервное правило.

\subsection{Порядок работы раунда}

Раунд устроен так:
\begin{enumerate}
\item бот отправляет ход МХ на сервер;
\item сервер возвращает ход соперника ХС;
\item по таблице~3 определяется результат РЗ;
\item результат учитывается в счётчиках СП (поражение) и СН (ничья);
\item запись о раунде добавляется в ИСТ (таблица~4);
\item если это был пятый раунд --- матч окончен, выводятся СП и СН;
\item иначе выбирается стратегия и вычисляется контрход КХ на следующий раунд;
\item $\text{МХ}:=\text{КХ}$, и начинается следующий раунд.
\end{enumerate}

Ход первого раунда фиксирован: $\text{МХ}:=\text{К}$, потому что истории ещё
нет и прогнозировать нечего.

% =====================================================================  5
\section{Таблицы, задающие автомат}

\begin{table}[H]
\centering
\caption{Таблица 1 --- контрход стратегии 1}
\label{tab:t1}
\begin{tabular}{|c|c|c|}
\hline
\textbf{ХС} & \textbf{КХ $=b(\text{ХС})$} & \textbf{Смысл}\\
\hline
К & Б & бумага бьёт камень\\
\hline
Н & К & камень бьёт ножницы\\
\hline
Б & Н & ножницы бьют бумагу\\
\hline
\end{tabular}
\end{table}

\begin{table}[H]
\centering
\caption{Таблица 2 --- контрход стратегии 2}
\label{tab:t2}
\begin{tabular}{|c|c|c|}
\hline
\textbf{СХС} & \textbf{КХ $=b(\text{СХС})$} & \textbf{Смысл}\\
\hline
К & Б & бот готовится к камню\\
\hline
Н & К & бот готовится к ножницам\\
\hline
Б & Н & бот готовится к бумаге\\
\hline
\end{tabular}
\end{table}

\begin{table}[H]
\centering
\caption{Таблица 3 --- определение результата РЗ}
\label{tab:t3}
\begin{tabular}{|c||c|c|c|}
\hline
\textbf{МХ $\backslash$ ХС} & \textbf{К} & \textbf{Н} & \textbf{Б}\\
\hline\hline
К & Н & В & П\\
\hline
Н & П & Н & В\\
\hline
Б & В & П & Н\\
\hline
\end{tabular}
\end{table}

В таблице~3: В --- бот выиграл раунд, Н --- ничья, П --- бот проиграл.
Диагональ --- ничьи.

\begin{table}[H]
\centering
\caption{Таблица 4 --- история ИСТ (пример заполнения для матча Н, К, Н, К, Н)}
\label{tab:t4}
\begin{tabular}{|c|c|c|c|}
\hline
\textbf{НР} & \textbf{ХС} & \textbf{МХ} & \textbf{РЗ}\\
\hline
1 & Н & К & В\\
\hline
2 & К & К & Н\\
\hline
3 & Н & Б & П\\
\hline
4 & К & Б & В\\
\hline
5 & Н & К & В\\
\hline
\end{tabular}
\end{table}

История пополняется одной записью за раунд; именно по ней работает
таблица~5.

\begin{table}[H]
\centering
\caption{Таблица 5 --- поиск паттерна (ИСТ $=$ Н, К, Н, К)}
\label{tab:t5}
\begin{tabular}{|c|c|c|c|}
\hline
\textbf{$k$} & \textbf{Паттерн} & \textbf{Более раннее вхождение} & \textbf{СХС}\\
\hline\hline
2 & (К, К) & --- & ---\\
\hline
2 & (К, Н) & --- & ---\\
\hline
2 & (К, Б) & --- & ---\\
\hline
2 & (Н, К) & раунды 1--2 & \textbf{Н}\\
\hline
2 & (Н, Н) & --- & ---\\
\hline
2 & (Н, Б) & --- & ---\\
\hline
2 & (Б, К) & --- & ---\\
\hline
2 & (Б, Н) & --- & ---\\
\hline
2 & (Б, Б) & --- & ---\\
\hline\hline
1 & (К) & раунд 2 & Н\\
\hline
1 & (Н) & раунд 1 & К\\
\hline
1 & (Б) & --- & ---\\
\hline
\end{tabular}
\end{table}

В примере из таблицы~5 последние два хода соперника --- (Н, К); такая пара уже
была в раундах 1--2, и после неё соперник сыграл Н. Значит $\text{СХС}=\text{Н}$
и бот отвечает по таблице~2 ходом К. Уровень $k=1$ в этом случае уже не нужен;
он включается, только когда пары ещё не повторялись (например, в третьем
раунде рассмотренного матча: пар нет, зато ход Н уже встречался и после него
был К, поэтому $\text{СХС}=\text{К}$).

% =====================================================================  6
\section{Схема алгоритма}

Схема (рис.~\ref{fig:flow}) построена по ГОСТ~19.701--90: овал --- начало и
конец, параллелограмм --- ввод/вывод, прямоугольник --- процесс, ромб ---
решение. Порядок блоков соответствует п.~4.5.

\begin{figure}[p]
\centering
\begin{tikzpicture}[
  x=1cm, y=1cm, font=\footnotesize,
  term/.style={draw, thick, ellipse, minimum width=3.2cm, minimum height=0.85cm,
               align=center, inner sep=1pt},
  proc/.style={draw, thick, rectangle, minimum width=5.4cm, minimum height=0.8cm,
               align=center, inner sep=1pt},
  io/.style={draw, thick, trapezium, trapezium left angle=72, trapezium right angle=108,
             minimum width=5.4cm, minimum height=0.8cm, align=center, inner sep=1pt},
  dec/.style={draw, thick, diamond, aspect=3.0, minimum width=3.4cm,
              minimum height=1.15cm, align=center, inner sep=0pt},
  wide/.style={dec, minimum width=4.4cm, aspect=3.4},
  arr/.style={-{Latex[length=1.8mm]}, thick},
  lbl/.style={font=\scriptsize, inner sep=1.5pt, fill=white}
]
% --- основной столбец
\node[term] (start) at (0,0) {\textbf{НАЧАЛО}};
\node[proc, minimum height=1.35cm] (init) at (0,-1.45)
      {Н.У.\\[1pt]НР := 0; \quad СП := 0; \quad СН := 0\\
       С := 0; \quad ИСТ := $\varnothing$};
\node[io] (srv1) at (0,-3.1) {Ответ сервера:\\ начало раунда};
\node[proc] (inc) at (0,-4.5) {НР := НР + 1};
\node[dec] (d1) at (0,-5.9) {НР = 1 ?};
\node[io] (io1) at (0,-7.5) {Вывод МХ\\ (ввод ХС --- ответ сервера)};
\node[proc] (rz) at (0,-9.0) {РЗ := табл.~3 (МХ, ХС)};
\node[dec] (dp) at (0,-10.4) {РЗ = П ?};
\node[dec] (dn) at (0,-11.8) {РЗ = Н ?};
\node[proc] (ist) at (0,-13.4) {ИСТ := ИСТ $\cup$ \{НР, ХС, МХ, РЗ\}};
\node[dec] (d5) at (0,-14.9) {НР = 5 ?};
\node[wide] (dsw) at (0,-17.3) {СП + СН $\ge$ 2 ?};
\node[proc, minimum width=4.6cm] (kx1) at (0,-19.1) {КХ := табл.~1 (по ХС)};
\node[proc, minimum width=4.0cm] (mx) at (0,-23.6) {МХ := КХ};
% --- правая колонка
\node[proc, minimum width=2.4cm] (mxk) at (4.6,-5.9) {МХ := К};
\node[proc, minimum width=2.6cm] (sp) at (4.6,-10.4) {СП := СП + 1};
\node[proc, minimum width=2.6cm] (sn) at (4.6,-11.8) {СН := СН + 1};
\node[io, minimum width=2.6cm] (out2) at (4.6,-14.9) {Вывод СП, СН};
\node[term, minimum width=2.2cm] (end) at (4.6,-16.3) {\textbf{КОНЕЦ}};
\node[proc, minimum width=2.2cm] (cs1) at (4.6,-17.3) {С := 1};
\node[proc, minimum width=2.9cm, minimum height=1.1cm] (search) at (4.6,-19.1)
      {Поиск паттерна в ИСТ\\ (табл.~5) $\to$ СХС};
\node[dec, minimum width=3.0cm] (dpatt) at (4.6,-20.9) {Паттерн найден ?};
\node[proc, minimum width=2.9cm] (kx2) at (4.6,-22.6) {КХ := табл.~2 (по СХС)};
\node[proc, minimum width=2.6cm] (kx1b) at (7.4,-20.9) {КХ := табл.~1 (по ХС)};

% --- поток
\draw[arr] (start) -- (init);
\draw[arr] (init) -- (srv1);
\draw[arr] (srv1) -- (inc);
\draw[arr] (inc) -- (d1);
\draw[arr] (d1.east) -- node[lbl,above]{да} (mxk.west);
\draw[arr] (d1.south) -- node[lbl,right]{нет} (io1.north);
\draw[arr] (mxk.south) -- (4.6,-7.5) -- (io1.east);
\draw[arr] (io1) -- (rz);
\draw[arr] (rz) -- (dp);
\draw[arr] (dp.east) -- node[lbl,above]{да} (sp.west);
\draw[arr] (dp.south) -- node[lbl,right]{нет} (dn.north);
\draw[arr] (dn.east) -- node[lbl,above]{да} (sn.west);
\draw[arr] (dn.south) -- node[lbl,right]{нет} (ist.north);
% слияние ветвей счётчиков
\draw[arr] (sp.east) -- (6.6,-10.4) -- (6.6,-13.4) -- (ist.east);
\draw[thick] (sn.east) -- (6.6,-11.8);
\fill (6.6,-11.8) circle (1.6pt);
\draw[arr] (ist) -- (d5);
\draw[arr] (d5.east) -- node[lbl,above]{да} (out2.west);
\draw[arr] (out2.south) -- (end.north);
\draw[arr] (d5.south) -- node[lbl,right]{нет} (dsw.north);
\draw[arr] (dsw.east) -- node[lbl,above]{да} (cs1.west);
\draw[arr] (dsw.south) -- node[lbl,right]{нет} (kx1.north);
\draw[arr] (cs1) -- (search);
\draw[arr] (search) -- (dpatt);
\draw[arr] (dpatt.south) -- node[lbl,right]{да} (kx2.north);
\draw[arr] (dpatt.east) -- node[lbl,above]{нет} (kx1b.west);
\draw[arr] (kx2.south) -- (4.6,-23.6) -- (mx.east);
\draw[thick] (kx1b.south) -- (7.4,-23.6) -- (4.6,-23.6);
\fill (4.6,-23.6) circle (1.6pt);
\draw[arr] (kx1) -- (mx);
% возврат к началу раунда
\draw[arr] (mx.west) -- (-4.6,-23.6) -- (-4.6,-4.5) -- (inc.west);
\end{tikzpicture}
\caption{Схема алгоритма работы бота (ГОСТ~19.701--90)}
\label{fig:flow}
\end{figure}

Пояснения к схеме:
\begin{itemize}
\item блок <<Н.У.>> задаёт~\eqref{eq:init}; никаких других начальных
      присваиваний в схеме нет;
\item <<Ответ сервера: начало раунда>> --- условный блок: турнир проводится на
      сервере, и раунд начинается только после его команды;
\item <<НР = 1 ?>> --- единственная проверка номера раунда: она нужна только
      для первого хода, дальше МХ всегда приходит из КХ;
\item <<РЗ = П ?>> и <<РЗ = Н ?>> --- две проверки результата, каждая
      увеличивает свой счётчик; обе ветви сходятся перед записью в ИСТ;
\item <<НР = 5 ?>> --- выход из цикла; при досрочном выходе КХ уже не
      вычисляется, лишних действий в последнем раунде нет;
\item <<СП + СН $\ge$ 2 ?>> и <<С := 1>> --- необратимое переключение на
      стратегию~2;
\item <<Паттерн найден ?>> --- развилка стратегии~2: при <<да>> работает
      таблица~2, при <<нет>> --- резервная таблица~1;
\item возврат в начало цикла идёт в блок <<НР := НР + 1>>, а не в начало схемы,
      поэтому Н.У. выполняется ровно один раз.
\end{itemize}

% =====================================================================  7
\section{Пример работы}

Рассмотрим матч, в котором соперник играет Н, К, Н, К, Н. Трассировка
приведена в табл.~\ref{tab:trace}.

\begin{table}[H]
\centering
\caption{Трассировка матча (соперник: Н, К, Н, К, Н)}
\label{tab:trace}
\begin{tabular}{|c|c|c|c|c|c|c|c|c|}
\hline
\textbf{НР} & \textbf{МХ} & \textbf{ХС} & \textbf{РЗ} & \textbf{СП} & \textbf{СН} &
\textbf{С} & \textbf{СХС} & \textbf{КХ}\\
\hline
1 & К & Н & В & 0 & 0 & 0 & --- & К\\
\hline
2 & К & К & Н & 0 & 1 & 0 & --- & Б\\
\hline
3 & Б & Н & П & 1 & 1 & 1 & К & Б\\
\hline
4 & Б & К & В & 1 & 1 & 1 & Н & К\\
\hline
5 & К & Н & В & 1 & 1 & 1 & --- & ---\\
\hline
\multicolumn{4}{|c|}{\textbf{Итог матча}} & \multicolumn{2}{c|}{В: 3 \quad Н: 1 \quad П: 1} &
\multicolumn{2}{c|}{} &\\
\hline
\end{tabular}
\end{table}

Разбор по раундам.

\begin{enumerate}
\item $\text{НР}=1$, поэтому $\text{МХ}:=\text{К}$. Соперник играет Н ---
      победа, $\text{СП}+\text{СН}=0<2$, стратегия остаётся первой:
      $\text{КХ}=b(\text{Н})=\text{К}$.
\item $\text{МХ}=\text{К}$, соперник --- К, ничья ($\text{СН}=1$). Сумма
      $\text{СП}+\text{СН}=1<2$: снова стратегия~1,
      $\text{КХ}=b(\text{К})=\text{Б}$.
\item $\text{МХ}=\text{Б}$, соперник --- Н, поражение ($\text{СП}=1$). Теперь
      $\text{СП}+\text{СН}=2$, поэтому $\text{С}:=1$. В ИСТ лежит (Н, К, Н):
      пары ещё не повторялись ($k=2$ --- нет), зато ход Н уже был в первом
      раунде, а после него соперник сыграл К. Значит $\text{СХС}=\text{К}$ и
      $\text{КХ}=b(\text{К})=\text{Б}$.
\item $\text{МХ}=\text{Б}$, соперник --- К, победа. В ИСТ уже (Н, К, Н, К):
      пара (Н, К) встречалась в раундах 1--2, после неё был Н, поэтому
      $\text{СХС}=\text{Н}$ и $\text{КХ}=b(\text{Н})=\text{К}$.
\item $\text{МХ}=\text{К}$, соперник --- Н, победа. $\text{НР}=5$ --- матч
      окончен, контрход не вычисляется.
\end{enumerate}

Итог: три победы, одна ничья, одно поражение. Оба прогноза (в третьем и
четвёртом раундах) оказались верными: соперник действительно играл по схеме
<<Н, К>> с повтором.

% =====================================================================  8
\section{Свойства алгоритма}

\subsection{Конечность}

Состояние бота полностью определяется историей ИСТ: номер раунда равен длине
истории, счётчики СП и СН --- это количества П и Н в ней, стратегия С равна
единице тогда и только тогда, когда в какой-то момент выполнялось
$\text{СП}+\text{СН}\ge2$, то есть тоже вычисляется по истории.

В ИСТ не может быть больше пяти записей, а каждый её элемент --- один из трёх
ходов. Число различных историй не превосходит
\[
1+3+9+27+81+243 = 364 .
\]
Значит, автомат конечен: у него не более 364 состояний, и он не может
<<раскачиваться>> бесконечно.

\subsection{Сложность}

\begin{itemize}
\item \textbf{Время.} Один раунд --- это фиксированное число действий: выбор
      хода по таблице, одна подстановка в таблицу~3, до двух сравнений
      результата, одна запись в историю и поиск паттерна. Поиск перебирает
      не более $3^2+3^1=12$ вариантов при фиксированной глубине, то есть
      выполняется за константу. Весь матч --- $O(5)$, то есть $O(1)$ при
      фиксированном числе раундов.
\item \textbf{Память.} Нужны история (не более пяти записей по два бита на ход
      плюс результат), два счётчика до пяти и один бит стратегии. Это
      несколько десятков бит --- на порядки меньше, чем у алгоритмов, хранящих
      полное дерево игры.
\item \textbf{Детерминированность.} Генератор случайных чисел не используется:
      при одной и той же последовательности ходов соперника бот всегда играет
      одну и ту же партию.
\end{itemize}

\subsection{Проверка перебором}

Программа из приложения~А перебирает все $3^5=243$ возможные
последовательности ходов соперника длиной пять. Результат:

\begin{itemize}
\item побед --- $39{,}5\,\%$ матчей, ничьих --- $21{,}0\,\%$,
      поражений --- $39{,}5\,\%$.
\end{itemize}

Симметрия чисел ожидаема: против \emph{равномерно случайного} соперника ни одна
детерминированная стратегия не может иметь перевеса в отдельном раунде.
Преимущество алгоритма проявляется на \emph{закономерных} соперниках. Для
периодических последовательностей суммарный баланс (победы минус поражения)
положителен: например, против постоянного хода --- $+4$ или $+5$ за пять
раундов, против периода 2 --- до $+4$.

Иными словами, алгоритм не портит игру против хаоса и выигрывает у того, кто
играет шаблонно, --- а именно шаблонных соперников и порождают типичные
студенческие реализации ботов.

% =====================================================================  9
\section{Выводы}

\begin{enumerate}
\item Построен детерминированный конечный автомат бота для турнира
      <<камень --- ножницы --- бумага>>, использующий ровно те переменные,
      которые заданы в Н.У.: МХ, ИСТ, НР, ХС, СХС, РЗ, КХ, СП, СН, С.
\item Метод состоит из двух стратегий: реактивной (контрход на фактический
      ход соперника) и проактивной (контрход на прогноз, полученный поиском
      паттерна в истории). Переключение происходит при
      $\text{СП}+\text{СН}\ge2$ и необратимо.
\item Работа автомата задана пятью таблицами: контрход стратегии~1, контрход
      стратегии~2, определение результата, история и поиск паттерна.
      Таблицы~1 и~2 используют одну функцию $b(\cdot)$ с разными аргументами,
      что объяснено в п.~4.3.
\item Схема алгоритма оформлена по ГОСТ~19.701--90 и начинается строго так:
      НАЧАЛО $\to$ Н.У. $\to$ ответ сервера $\to$ алгоритм. Условный блок
      <<Паттерн найден ?>> расположен в конце цикла, перед выбором контрхода.
\item Конечность автомата доказана оценкой числа историй (не более 364),
      работа проверена трассировкой матча и полным перебором всех 243 партий.
\end{enumerate}

% =====================================================================  ЛИТЕРАТУРА
\section*{Список литературы}
\addcontentsline{toc}{section}{Список литературы}
\begin{enumerate}
\item ГОСТ 19.701--90. Единая система программной документации. Схемы
      алгоритмов, программ, данных и систем. Обозначения условные и правила
      выполнения. --- М.: Изд-во стандартов, 1991.
\item Кудрявцев В.~Б., Алешин С.~В., Подколзин А.~С. Введение в теорию
      автоматов. --- М.: Наука, 1985.
\item Карпов Ю.~Г. Теория автоматов. --- СПб.: Питер, 2003.
\end{enumerate}

% =====================================================================  ПРИЛОЖЕНИЕ
\newpage
\appendix
\section{Листинг программы проверки}

{\footnotesize
\begin{verbatim}
# -*- coding: utf-8 -*-
MOVES = ['К', 'Н', 'Б']
BEAT = {'К': 'Б', 'Н': 'К', 'Б': 'Н'}    # контрход
BEATS = {'Б': 'К', 'К': 'Н', 'Н': 'Б'}   # что бьёт ход

N_ROUNDS = 5
SWITCH = 2                               # порог СП + СН >= 2

def result(my, opp):                     # таблица 3
    if my == opp:
        return 'Н'
    return 'В' if BEATS[my] == opp else 'П'

T3 = {m: {o: result(m, o) for o in MOVES} for m in MOVES}
T1 = {o: BEAT[o] for o in MOVES}         # таблица 1
T2 = {o: BEAT[o] for o in MOVES}         # таблица 2

def search_pattern(hist, ks=(2, 1)):     # таблица 5
    for k in ks:
        if len(hist) < k + 1:
            continue
        pat = tuple(r['ХС'] for r in hist[-k:])
        found = None
        for i in range(len(hist) - k - 1):
            if tuple(r['ХС'] for r in hist[i:i + k]) == pat:
                found = hist[i + k]['ХС']
        if found is not None:
            return found, k
    return None, None

def play(opp_seq):
    НР, СП, СН, С = 0, 0, 0, 0           # Н.У.
    ИСТ, МХ = [], None
    rows = []
    for t, ХС in enumerate(opp_seq, 1):
        НР = t
        if НР == 1:
            МХ = 'К'
        РЗ = T3[МХ][ХС]
        if РЗ == 'П':
            СП += 1
        elif РЗ == 'Н':
            СН += 1
        ИСТ.append({'НР': НР, 'ХС': ХС, 'МХ': МХ, 'РЗ': РЗ})
        row = {'НР': НР, 'МХ': МХ, 'ХС': ХС, 'РЗ': РЗ,
               'СП': СП, 'СН': СН, 'С': С, 'СХС': None,
               'паттерн': False, 'КХ': None}
        if НР == N_ROUNDS:
            rows.append(row)
            break
        if СП + СН >= SWITCH:
            С = 1
        row['С'] = С
        if С == 0:
            КХ = T1[ХС]
        else:
            СХС, lvl = search_pattern(ИСТ)
            if СХС is not None:
                row['паттерн'] = True
                row['СХС'] = СХС
                КХ = T2[СХС]
            else:
                КХ = T1[ХС]
        row['КХ'] = КХ
        rows.append(row)
        МХ = КХ
    return rows

def score(rows):
    return (sum(r['РЗ'] == 'В' for r in rows),
            sum(r['РЗ'] == 'Н' for r in rows),
            sum(r['РЗ'] == 'П' for r in rows))

if __name__ == '__main__':
    for r in play(['Н', 'К', 'Н', 'К', 'Н']):
        print(r['НР'], r['МХ'], r['ХС'], r['РЗ'], r['СП'], r['СН'],
              r['С'], r['СХС'], r['КХ'])

    from itertools import product
    W = D = L = 0
    for seq in product(MOVES, repeat=N_ROUNDS):
        v, n, p = score(play(list(seq)))
        W += v > p
        D += v == p
        L += v < p
    print('матчей:', W + D + L, ' В:', W, ' Н:', D, ' П:', L)
\end{verbatim}
}

\end{document}
